\subsection{Выбор используемого метода по обнаружению аномалий в контейнерных средах}

На основе проведённого анализа существующих подходов и методов обнаружения аномалий в распределённых и контейнерных средах можно сделать вывод о том, что наиболее перспективными являются методы, основанные на нейронных сетях. Это обусловлено тем, что такие методы способны выявлять сложные нелинейные зависимости в данных, учитывать временные и контекстные связи между событиями, а также адаптироваться к изменяющемуся поведению системы.

Ключевой особенностью рассматриваемой задачи является отсутствие размеченных данных в процессе обучения. В реальных Kubernetes-кластерах аномальные события являются редкими, неполными или заранее неизвестными, что делает невозможным формирование полноценного обучающего набора с корректной разметкой классов. В таких условиях применение методов контролируемого обучения (supervised learning) существенно ограничено и не обеспечивает требуемой обобщающей способности модели.

Дополнительным усложняющим фактором является то, что каждый Kubernetes-кластер обладает индивидуальной архитектурой, набором развернутых сервисов, политиками безопасности, конфигурацией сетевого взаимодействия и характером нагрузки. В результате поведение «нормальной» системы существенно варьируется от кластера к кластеру, а также изменяется во времени в рамках одного и того же кластера. Это приводит к тому, что универсальный набор шаблонов или правил поведения практически невозможно сформировать без существенных ограничений по применимости.

В связи с этим в рамках данной работы рассматриваются методы неконтролируемого обучения (unsupervised learning), позволяющие обучать модель исключительно на данных нормального поведения системы. К числу таких методов относятся нейронные сети на основе автоэнкодеров (Autoencoder, AE).

В качестве основных рассматриваемых архитектур выбраны следующие варианты автоэнкодеров:

\begin{itemize}
    \item \textbf{Автоэнкодер (AE)} — базовая архитектура, состоящая из энкодера и декодера, предназначенная для восстановления входных данных и выявления аномалий по величине ошибки реконструкции;
    \item \textbf{RNN-AE} — автоэнкодер, в котором энкодер и декодер реализованы с использованием рекуррентных нейронных сетей, что позволяет учитывать последовательностную структуру входных данных;
    \item \textbf{LSTM-AE} — расширение RNN-AE с использованием LSTM-слоёв, обеспечивающих более эффективное моделирование долгосрочных зависимостей в последовательностях логов и событий.
\end{itemize}

Выбор данных архитектур обусловлен их способностью работать с неразмеченными данными, характерными для задач мониторинга контейнерных сред, а также их доказанной эффективностью в задачах обнаружения аномалий на основе ошибки реконструкции. В рамках дальнейших разделов работы будет проведено проектирование и реализация выбранных моделей, а также их сравнительный анализ в контексте применения для обнаружения аномалий в Kubernetes-кластере.